\bigskip
\noindent Autres pistes de l'onglet Findings pour ce thème, dans des cotes qui n'ont pas de section ici :\par

\begin{encadre}
\begin{theoreme*}[Quillen, d'après les pages 43 à 45, universalité des correspondances de bordisme\piste{162-5-bordism-correspondences-universal}]\nonverifie
Sur les variétés fermées et les classes d'homotopie d'applications, la catégorie $\mathbf B$ dont les morphismes $X \to Y$ sont les classes de bordisme non orienté de variétés au-dessus de $X \times Y$ est universelle parmi les catégories munies d'un foncteur covariant et d'un foncteur contravariant qui coïncident sur les objets et vérifient le changement de base pour les carrés cartésiens transverses. Tout tel couple $(F'_\bullet, F'^\bullet)$ identifie les correspondances bordantes, une correspondance $Z \to X \times Y$ de composantes $p$, $q$ allant sur $F'_\bullet(q) \circ F'^\bullet(p)$. La page attribue l'énoncé à Quillen.
\end{theoreme*}
\end{encadre}
{\small Les feuillets en esquissent une démonstration incomplète (pages 43 à 45). L'esquisse finit par \q{Le reste est sans doute l'AQT}, sans rien sur la composition ni sur l'unicité, et le mot \q{fermées} est de l'édition. Mathlib n'a pas le bordisme des variétés. On se reportera aux notes de Grothendieck (\lot{162-5}{3}{43}).\par}

\begin{encadre}
\begin{theoreme*}[Grothendieck, pages 27 à 30, de Rham et Hochschild pour une catégorie formelle\piste{162-5-de-rham-hochschild-flat-formal-category}]\nonverifie
Soient $X_0$ un schéma de caractéristique nulle et $\mathcal C$ une catégorie formelle sur $X_0$ de conormal $\Omega$ plat. Les données de descente sur un $\mathcal O$-module $M$ relativement à $\mathcal C$ correspondent aux $\Omega$-connexions intégrables sur $M$, et le complexe de de Rham $\Lambda^\bullet\Omega \otimes M$ est canoniquement quasi-isomorphe, dans la catégorie dérivée, au complexe de Hochschild $C^\bullet(\mathcal C, M)$.
\end{theoreme*}
\end{encadre}
{\small La page~30 l'énonce sans démonstration, \q{sous les conditions du th. de Quillen} ; le mot \q{intégrables} et la définition de $C^\bullet(\mathcal C, M)$ sont de la lecture. Mathlib n'a pas les catégories formelles ni les complexes de de Rham sur un schéma. On se reportera aux notes de Grothendieck (\lot{162-5}{2}{27}).\par}

\begin{encadre}
\begin{theoreme*}[Quillen, d'après les pages 21 à 26, catégories formelles à conormal plat\piste{162-5-formal-categories-flat-conormal-de-rham}]\nonverifie
Sur un schéma $X_0$ de caractéristique nulle, le foncteur qui envoie une catégorie formelle $I$-adique sur $X_0$ (un pro-anneau complet augmenté muni d'un $\Delta$ coassociatif et counitaire, sans inverse supposé) sur l'algèbre différentielle graduée $(\Lambda^\bullet\Omega, \delta)$ de son conormal $\Omega = I/I^2$ est une équivalence entre les catégories formelles de $\Omega$ plat et les complexes de de Rham sur $\Lambda^\bullet\Omega$ de $\Omega$ plat, sans hypothèse de finitude sur $\Omega$ ni de lissité sur~$X_0$.
\end{theoreme*}
\end{encadre}
{\small La page~26 l'énonce comme \q{Th. de Quillen}, sans démonstration ni référence, et dit qu'il contient le théorème de Cartan sur les groupes formels en caractéristique nulle. Mathlib n'a pas les catégories formelles $I$-adiques. On se reportera aux notes de Grothendieck (\lot{162-5}{2}{21}).\par}
